\chapter{中断、设备与用户边界}

\section{从轮询到中断}

轮询由 CPU 主动反复读取状态；中断由设备在状态变化时请求处理。启用中断前，
内核必须完成 IDT、控制器屏蔽、EOI 规则和共享状态同步。中断上半部只确认和
记录事件，较长工作应延后执行。

\section{PIC、PIT、键盘与 IDE}

PIC 重映射避免 IRQ 与 CPU 异常向量冲突；PIT 提供周期时钟；键盘控制器提供
扫描码；IDE 提供扇区 I/O。每个驱动都应分为寄存器访问、状态机、策略和公共
接口，防止端口魔法数字散落在调度器或文件系统中。

\section{Ring 3 与系统调用}

用户态不能直接访问内核页和特权指令。进入 Ring 3 需要用户代码/数据段、
用户栈、页表权限和正确的返回帧。系统调用入口必须先保存不可信上下文，再验证
编号、指针、长度和权限；验证通过才触碰内核对象。

\begin{contractbox}
用户指针不是 C++ 指针承诺，而是“可能映射、可能跨页、可能在验证后变化”的
不可信地址。复制进出用户空间必须定义部分失败和页故障处理。
\end{contractbox}

\section{驱动的完成语义}

命令写入设备不等于操作完成。驱动必须区分提交、设备接受、数据可用、完成确认
和错误恢复。这个模型会一直延伸到块层、文件系统和用户 \code{read} 返回值。
